\section{Chapter~\ref{sec:acet}. Adding \nt{ACET}}

The three actions of acetylcholine are measured in slices and in the living brain, the conflict between writing and reading is shown with the book's model, and the idea that \nt{ACET} serves as a write-enable line and reports expected uncertainty is a hypothesis with partial experimental support.

{\footnotesize
\setlength{\tabcolsep}{3pt}\renewcommand{\arraystretch}{1.15}
\begin{longtable}{>{\raggedright}p{0.5cm}>{\raggedright}p{4.0cm}>{\raggedright}p{5.6cm}>{\raggedright}p{2.3cm}>{\raggedright\arraybackslash}p{1.7cm}}
\toprule
No. & Claim & Experiment and what was measured & Source & Status \\
\midrule\endfirsthead
\toprule
No. & Claim & Experiment and what was measured & Source & Status \\
\midrule\endhead
1 & The brain's \nt{ACET} neurons are few, gathered in a few groups, and their outputs spread throughout the cortex: a broadcast channel & Antibody staining for the acetylcholine-synthesizing enzyme and retrograde labeling in the rat: the cortex, hippocampus, amygdala, olfactory bulb and thalamus receive acetylcholine mainly from six compact cell groups of the basal forebrain and brainstem & \cite{Mesulam1983} & measured \\
2 & The acetylcholine level in the cortex is high in waking and low in slow-wave sleep & Microdialysis in the cortex and hippocampus of freely moving cats with EEG recording: acetylcholine release is $\sim100\%$ higher in quiet waking and $\sim175\%$ higher in active waking than in slow-wave sleep; in REM sleep, cortical release is as in waking & \cite{Marrosu1995,Hasselmo1999} & measured \\
3 & The meaning of the signal is set by the receptor: acetylcholine has fast channel receptors and slow ones that trigger reactions in the cell (Proposition~\ref{prop:receptor}) & Review of measurements: the action of acetylcholine on excitability, transmission and plasticity depends on the release site, the receptor subtype (nicotinic are ion channels, muscarinic act through G proteins) and the recipient cell & \cite{Picciotto2012} & measured \\
4 & First action: weaken the internal connections while hardly touching external inputs (factor $1-a$ in~\eqref{eq:acetnet}) & Rat olfactory cortex slices: at $100$\,\textmu M acetylcholine agonists, potentials of internal connections (layer Ib) fall by more than $60\%$, those of the external input (layer Ia) by less than $15\%$, in the same cell; the suppression is presynaptic. Hippocampal slices (CA1): $89\%$ versus $40\%$ for the internal and input pathways & \cite{HasselmoBower1992,HasselmoSchnell1994} & measured \\
5 & Second action: strengthen the input (factor $\frac{1+a}{2}$) & Neocortical slices: nicotinic receptors strengthen synapses of the thalamic input and do not act on intracortical ones; muscarinic receptors weaken both types. Acetylcholine raises neuronal excitability (review) & \cite{Gil1997,Hasselmo2006} & measured \\
6 & Third action: make weight changes easier (rate $\eta a$) & Rat: electrical stimulation of the nucleus basalis (the source of cortical acetylcholine) paired with a tone produces, in an adult animal, a strong reorganization of the auditory cortex map toward the presented sound & \cite{KilgardMerzenich1998,Hasselmo2006} & measured \\
7 & The Hebb rule for sparse patterns in the second equation of~\eqref{eq:acetnet} gives an associative memory that completes a pattern from a part & Capacity calculation for a network with sparse patterns and a covariance rule: at a small fraction $p$ of active neurons, the network stores considerably more patterns than at $p=1/2$ & \cite{Tsodyks1988} & theory \\
8 & Writing at low $a$ corrupts memory: the network writes down what it recalled, not what it saw; at $a=0.1$ the corruption is no weaker than at $a=0$ & \texttt{models/\allowbreak acet\_memory.py}: $200$ neurons, five patterns of $20$, a new pattern shares $10$ neurons with one of them, $10$ runs. At $a=0$ the new pattern is not memorized, and the old one drags along $5.9$ foreign neurons out of $10$; at $a=0.1$ the new pattern is recalled ($0.97$), but the two merge: the new one drags along $7.3$ foreign neurons, the old one $7.9$; from $a=0.2$ on, fewer than one & derived in the chapter & model \\
9 & Proposition~\ref{prop:writeread}: there is no level $a$ at which writing and reading are equally good (Fig.~\ref{fig:acetsim}) & The same script, learning rate $\eta a$: a new pattern is memorized in full after one presentation for $a\ge0.9$, to $0.8$ at $a=0.75$, not at all for $a\le0.5$. Reading from half a pattern: $1.0$ for $a\le0.2$, $0.96$ at $a=0.3$, $0.04$ at $a=0.4$ & derived in the chapter & model \\
10 & In the brain a single broadcast wire, \nt{ACET}, switches between writing and reading & The idea comes from models built on slice measurements. The most direct experiment is in humans: the muscarinic blocker scopolamine impairs learning of new word pairs, most of all those overlapping with old ones, but not recall of pairs already learned. There is no direct measurement of two network modes & \cite{HasselmoAndersonBower1992,Atri2004} & hypothesis \\
11 & The filter~\eqref{eq:kalman} is optimal; the input weight and the learning rate are one quantity, $P/R$; in steady state $P=\sqrt{qR}$ and the memory is $\sqrt{R/q}$ (Proposition~\ref{prop:kalman}) & Requires no experiment: the optimality of the Kalman--Bucy filter is a classical result of estimation theory; the steady state is derived in the chapter & \cite{KalmanBucy1961} & theory \\
12 & \nt{ACET} tells the network a quantity like $P$: the expected uncertainty & Proposed in a probabilistic model of attention. There is no direct measurement showing that the acetylcholine level tracks the uncertainty of an estimate & \cite{YuDayan2005} & hypothesis \\
13 & Some \nt{ACET} neurons respond to reward and punishment within twenty milliseconds or so, more strongly the more unexpected the reinforcement & Mouse, optogenetically identified cholinergic neurons of the basal forebrain: response to reward and punishment after $18\pm3$\,ms; the size grows with the surprise of the reinforcement; responses are similar in neurons of two nuclei & \cite{Hangya2015} & measured \\
14 & \nt{NORA} notices an unexpected change in the world, and \nt{ACET} keeps the network in write mode & The division of labor is proposed in a model. Indirectly: in humans, pupil diameter, linked to \nt{NORA}, tracks changes and the learning rate. There is no direct test of the \nt{NORA}--\nt{ACET} pair & \cite{YuDayan2005,Nassar2012} & hypothesis \\
15 & In slow-wave sleep the network, in read mode, replays what was written during the day, and these replays rewrite it into long-term memory & Recordings of rat hippocampal ensembles: cells that worked together during the day fire together more often in subsequent slow-wave sleep, and in the same order (measured). For rewriting: suppressing hippocampal ripples after learning impairs spatial memory, and lengthening spontaneous ripples improves it; that the cause is a low \nt{ACET} level has not been proven & \cite{WilsonMcNaughton1994,Skaggs1996,Girardeau2009,FernandezRuiz2019,Hasselmo1999} & hypothesis \\
\bottomrule
\end{longtable}}
